\emph{ટ્યુરિંગ મશીન} સંગણનની આદર્શીકૃત યંત્રણા છે. તેમાં એક તરફ અનંત
\emph{પટ્ટી} હોય છે. પટ્ટી ખાનાંમાં વહેંચાયેલી છે. દરેક ખાનામાં અગાઉથી
નક્કી કરેલા વર્ણમાળાગણનું એક \emph{ચિહ્ન} રહી શકે છે. પટ્ટી પર
\emph{વાંચન-લેખન શીર્ષ} ખસેડીને મશીન કાર્ય કરે છે. મશીન અગાઉથી નક્કી કરેલી
\emph{અવસ્થાઓ} પૈકી એકમાં હોય છે. સૂચનાઓનો ગણ શીર્ષની ક્રિયા નક્કી કરે છે.
દરેક સૂચના મશીનની હાલની અવસ્થા અને તે વાંચે છે તે ચિહ્ન પર આધારિત છે.
હાલના ખાનામાં શું લખવું, શીર્ષને એક ખાનું ડાબે કે જમણે ખસેડવું કે સ્થિર
રાખવું, અને કઈ અવસ્થામાં જવું, તે સૂચના નિર્દિષ્ટ કરે છે. પટ્ટી પર
ચિહ્નોના ક્રમરૂપે \emph{નિવેશ} મૂક્યો હોય અને મશીનને આરંભિક અવસ્થામાં
રાખી શીર્ષ નિવેશના સૌથી ડાબા ખાનાને વાંચતું હોય, તો સૂચનાઓ મશીનના
\emph{સંરૂપણો}નો ક્રમ પગલે પગલે નક્કી કરે છે. દરેક સંરૂપણ પટ્ટીની
સામગ્રી, શીર્ષનું સ્થાન અને મશીનની અવસ્થા આપે છે. જે સંરૂપણમાં હાલના
ચિહ્ન અને અવસ્થા માટે સૂચના જ ન હોય ત્યાં મશીન \emph{અટકે} છે અને પટ્ટી
પરની સામગ્રી \emph{નિર્ગમ} છે.

ટ્યુરિંગ મશીનની પટ્ટી પર સંખ્યાઓને ઊભી લીટીઓના ક્રમથી સરળતાથી રજૂ કરી
શકાય છે. $\Nat \to \Nat$ પ્રકારનું ફલન \emph{ટ્યુરિંગ-સંગણનીય} છે
જો એવું ટ્યુરિંગ મશીન હોય કે જેને $n$ની એકાધારી રજૂઆત નિવેશ તરીકે આપતાં
તે અંતે અટકે અને પટ્ટી પર $f(n)$ની એકાધારી રજૂઆત નિર્ગમ તરીકે આપે.
ઘણાં પરિચિત અંકગણિતીય ફલનો ટ્યુરિંગ-સંગણનીય છે. કેટલાંક માટે આ
બતાવવું સરળ છે અને કેટલાંક માટે અઘરું. ટ્યુરિંગ મશીન સિવાય સંગણનના બીજા
ઘણા નિદર્શો પ્રસ્તાવિત થયા છે. એમનામાં સંગણનીય એવાં અંકગણિતીય ફલનો
દરેક વખતે ટ્યુરિંગ-સંગણનીય પણ નીકળ્યાં છે. આને \emph{ચર્ચ--ટ્યુરિંગ
પ્રતિજ્ઞા}નું સમર્થન ગણાય છે. એ પ્રતિજ્ઞા મુજબ અસરકારક રીતે ગણી શકાય
એવું દરેક અંકગણિતીય ફલન ટ્યુરિંગ-સંગણનીય છે.
